\clearpage

\section {Introduction }
\label{sec:intro}

In this technical note, 
we report the results of a 2-flavour $\nu_{\mu}$-disappearance analysis,
performed by the Rutherford-Oxford-Lancs groups,	
on the combined Run1+2 (1.431$\times$10$^{20}$ POT) dataset.
This dataset was collected in two run periods:
in January 2010 - June 2010 (Run1, 0.323$\times$10$^{20}$ POT) and
in November 2010 - March 2011 (Run2, 1.108$\times$10$^{20}$ POT).

This analysis 
\footnote{
   Preliminary results from this work were presented on May 20th 2011
   at the T2K collaboration meeting at KEK \cite{OurPreliminaryResults}.
   Sensitivity studies, performed before looking at the Run1+2 dataset,
   were presented on March 31st 2011 \cite{OurSensitivityStudy}.
}
is an update (with larger statistics) of the work 
that was reported in Ref.\ \cite{OurRun1DisappAnalysis},
and reviewed by M.Nakahata et al.\ \cite{Run1DisappAnalysisReview} in February 2011.
This note is meant to be read in conjunction with Ref.\ \cite{OurRun1DisappAnalysis} as the methodology is largely unchanged.
The method for predicting the nominal SuperK reconstructed energy spectrum of 1 $\mu$-like ring events is presented 
in Sections 2.1 - 2.6 of Ref.\ \cite{OurRun1DisappAnalysis} and is unchanged.
This analysis uses an updated SuperK flux tuning and an updated 
ND280 event-rate normalisation correction compared with those used to obtain the results in Ref.\ \cite{OurRun1DisappAnalysis}.
The flux normalisation ($N_{sk}/N_{nd}$) error, the flux shape error, the SuperK energy scale error and the 
SuperK ring-counting efficiency error have also been revised.
All changes to the analysis inputs are summarised in Sec.\ \ref{sec:analysis_modifications} of this note.
Updated nominal 1.431$\times$10$^{20}$ POT 1 $\mu$-like ring event predictions are shown in 
Sec. \ref{sec:updated_predictions} of this note.

As in the previous analysis, the oscillation fit was made using the likelihood ratio method \cite{StatAnalysisCowan}, 
and the confidence regions were constructed using the Feldman-Cousins method \cite{FeldmanCousins}; 
these are unchanged and are described in 
Sections 3.1, 3.2 and 3.4 of Ref.\ \cite{OurRun1DisappAnalysis}.
Systematic errors were included in the oscillation contours 
using the Cousins-Highland method (Bayesian integration) and by determining their effects on the 
Feldman-Cousins $\Delta\chi^{2}_{critical}$ values as in Ref.\ \cite{OurRun1DisappAnalysis}.
The $\Delta\chi^{2}_{critical}$ values have been updated compared with Ref. \cite{OurRun1DisappAnalysis}
using the modified analysis inputs. This updated Feldman-Cousins construction is discussed in 
Sec.\ \ref{sec:fc} of this note.

The results of the measurement of sin$^{2}$2$\theta_{23}$ and $\Delta m^{2}_{23}$ 
using the Run1+2 dataset are given in Sec. \ref{sec:fits} of this note 
for both GENIE and NEUT. The Feldman-Cousins confidence regions are shown in Sec.\ \ref{sec:conf_regions}.
The results from the fit to the combined Run1+2 dataset are compared with external measurements,
the T2K sensitivity (for the MINOS 2011 best-fit oscillation parameters) 
and the results of separate fits to the Run1-only and Run2-only datasets.

Finally, a summary is given in section \ref{sec:summary}.  

